% 图 4.4 性能对比柱状图（重绘加强版）
% 基于 CICC0903540 技术文档 4.2.2 - TVM 优化效果
\documentclass[tikz,border=10pt]{standalone}

% 强制全局样式与配色
\usepackage{amssymb, amsmath}
\usepackage{fontspec}
\usepackage{xeCJK}
\setCJKmainfont{SimHei}
\usepackage{tikz}
\usepackage{pgfplots}
\pgfplotsset{compat=1.18}
\usetikzlibrary{
  arrows.meta, positioning, shapes.geometric, calc,
  backgrounds, fit, decorations.pathreplacing, patterns
}

% 低饱和度马卡龙配色
\definecolor{mBlue}{HTML}{AEC6CF}
\definecolor{mPink}{HTML}{FFD1DC}
\definecolor{mGreen}{HTML}{B1D8B7}
\definecolor{mPurple}{HTML}{B39EB5}
\definecolor{mYellow}{HTML}{FDFD96}
\definecolor{mGray}{HTML}{E5E4E2}
\definecolor{mDark}{HTML}{4A4A4A}

% 全局 TikZ 样式
\tikzset{
  >=stealth,
  thick,
  draw=mDark,
  legendbox/.style={
    rectangle, draw=mDark!50, rounded corners=3pt,
    fill=mGray!15, inner sep=5pt, font=\tiny\sffamily},
}

\begin{document}
\begin{tikzpicture}
  \begin{axis}[
    ybar,
    bar width=0.8cm,
    width=14cm,
    height=8cm,
    ylabel={推理时间 (ms/image)},
    ylabel style={font=\sffamily\small, color=mDark},
    xlabel={优化方案对比（Places365 数据集，300张图片）},
    xlabel style={font=\sffamily\small, color=mDark},
    symbolic x coords={
      Baseline\\(Payload),
      Current\\(Payload),
      Rebuild\\(Current),
      Incremental\\(Current),
      Baseline\\(端到端),
      Current\\(端到端),
      MNN\\(端到端)
    },
    xtick=data,
    x tick label style={font=\sffamily\scriptsize, color=mDark, align=center},
    ymin=0, ymax=2600,
    ytick={0, 500, 1000, 1500, 2000, 2500},
    grid=major,
    grid style={draw=mGray, thin},
    legend style={
      at={(0.5,1.05)},
      anchor=south,
      legend columns=3,
      font=\sffamily\tiny,
      draw=mDark!50,
      fill=mGray!15,
      rounded corners=2pt
    },
    nodes near coords,
    nodes near coords style={
      font=\sffamily\tiny,
      color=mDark,
      /pgf/number format/.cd,
      fixed,
      precision=1,
      /pgf/number format/assume math mode
    },
  ]
    % Payload 对比（Scheme A）
    \addplot[fill=mPink!60, draw=mDark] coordinates {
      (Baseline\\(Payload), 1829.28)
      (Current\\(Payload), 0)
      (Rebuild\\(Current), 0)
      (Incremental\\(Current), 0)
      (Baseline\\(端到端), 0)
      (Current\\(端到端), 0)
      (MNN\\(端到端), 0)
    };

    \addplot[fill=mGreen!60, draw=mDark] coordinates {
      (Baseline\\(Payload), 0)
      (Current\\(Payload), 152.846)
      (Rebuild\\(Current), 0)
      (Incremental\\(Current), 0)
      (Baseline\\(端到端), 0)
      (Current\\(端到端), 0)
      (MNN\\(端到端), 0)
    };

    % Current-only 对比（Scheme B）
    \addplot[fill=mYellow!60, draw=mDark] coordinates {
      (Baseline\\(Payload), 0)
      (Current\\(Payload), 0)
      (Rebuild\\(Current), 2479.246)
      (Incremental\\(Current), 0)
      (Baseline\\(端到端), 0)
      (Current\\(端到端), 0)
      (MNN\\(端到端), 0)
    };

    \addplot[fill=mBlue!60, draw=mDark] coordinates {
      (Baseline\\(Payload), 0)
      (Current\\(Payload), 0)
      (Rebuild\\(Current), 0)
      (Incremental\\(Current), 152.36)
      (Baseline\\(端到端), 0)
      (Current\\(端到端), 0)
      (MNN\\(端到端), 0)
    };

    % 真实端到端重建（最终结论）
    \addplot[fill=mPink!60, draw=mDark] coordinates {
      (Baseline\\(Payload), 0)
      (Current\\(Payload), 0)
      (Rebuild\\(Current), 0)
      (Incremental\\(Current), 0)
      (Baseline\\(端到端), 1830.3)
      (Current\\(端到端), 0)
      (MNN\\(端到端), 0)
    };

    \addplot[fill=mPurple!60, draw=mDark] coordinates {
      (Baseline\\(Payload), 0)
      (Current\\(Payload), 0)
      (Rebuild\\(Current), 0)
      (Incremental\\(Current), 0)
      (Baseline\\(端到端), 0)
      (Current\\(端到端), 255.931)
      (MNN\\(端到端), 0)
    };

    \addplot[fill=mGreen!60, draw=mDark] coordinates {
      (Baseline\\(Payload), 0)
      (Current\\(Payload), 0)
      (Rebuild\\(Current), 0)
      (Incremental\\(Current), 0)
      (Baseline\\(端到端), 0)
      (Current\\(端到端), 0)
      (MNN\\(端到端), 410)
    };

    \legend{
      Baseline (Payload),
      Current (Payload),
      Rebuild-only,
      Incremental,
      Baseline (真实),
      Current (真实),
      MNN
    }
  \end{axis}

  % ========== 信息密度拉满：性能提升标注 ==========
  \node[legendbox, anchor=north west] at (0.5, 7.5)
    [align=left, text width=5.5cm] {
    \textbf{性能提升归因分析}\\[2pt]
    • Payload 对比：91.64\% 提升（1829 $\to$ 153 ms）\\
    • Current-only：93.85\% 提升（2479 $\to$ 152 ms）\\
    • 端到端重建：86.02\% 提升（1830 $\to$ 256 ms）\\
    • MNN 动态支持：1.85$\times$ 加速（平均 410 ms）
  };

  % ========== 图例（右下角）==========
  \node[legendbox, anchor=south east] at (13.5, 0.5)
    [align=left, text width=4.8cm] {
    \textbf{图例}\\[1pt]
    \colorbox{mPink!60}{\rule{0.16cm}{0.12cm}\hspace{0.05cm}} Baseline 路线\\
    \colorbox{mGreen!60}{\rule{0.16cm}{0.12cm}\hspace{0.05cm}} Current 路线\\
    \colorbox{mPurple!60}{\rule{0.16cm}{0.12cm}\hspace{0.05cm}} TVM 端到端\\
    \colorbox{mYellow!60}{\rule{0.16cm}{0.12cm}\hspace{0.05cm}} Rebuild-only\\[1pt]
    测试平台：飞腾派 ARMv8.0 Cortex-A72\\
    测试数据：Places365（256$\times$256, 300张）
  };
\end{tikzpicture}
\end{document}
